\documentclass[10pt,a4paper]{amsart}
\usepackage[margin=19mm]{geometry}
\usepackage{amsmath,amssymb,mathtools}
\usepackage[scr=rsfs,cal=euler]{mathalfa}
\usepackage{xfrac}
\usepackage[unicode,hidelinks,bookmarks=false]{hyperref}
\newcommand{\ie}{i.e.\ }
\newcommand{\cf}{cf.\ }
\newcommand{\resp}{respectively\ }
\newcommand{\Subsection}{\subsection}
\providecommand{\nobreakdash}{\nobreak\hbox{-}}
\setlength{\emergencystretch}{2em}
\multlinegap0pt
\usepackage{amssymb}
\usepackage[shortlabels]{enumitem}
\usepackage{mathtools}
\newtheorem{lem}{Lemma}
\newtheorem{cor}{Corollary}
\newcommand{\B}{\mathcal{B}}%{\mathscr{B}}
\newcommand{\F}{\mathcal{F}}
\newcommand{\R}{{\mathbb{R}}}
\title{Statistical Inference for Quasi-Infinitely Divisible Distributions via Fourier Methods}
\author{Vladimir Panov, Anton Ryabchenko}
\date{Source: arXiv:2505.14255v2 (CC BY 4.0)}
\begin{document}
\maketitle

\noindent\emph{Source and licence.} Statistical Inference for Quasi-Infinitely Divisible Distributions via Fourier Methods. Version arXiv:2505.14255v2 (CC BY 4.0), distributed by arXiv under the Creative Commons Attribution 4.0 International licence, http://creativecommons.org/licenses/by/4.0/, as recorded in the arXiv licence metadata for this exact version and shown on the abstract page https://arxiv.org/abs/2505.14255v2. Full licence notice: \texttt{licenses/arxiv-2505.14255-CC-BY-4.0.txt}.\par\smallskip

\noindent\textit{Editorial scope.} Counted unit: the article's cor cor (source0.tex lines 375--377). The article's complete original proof of it is delivered with it (source0.tex lines 378--387, one \texttt{proof} environment of 10 lines). No mathematics is written, repaired or re-derived: the statement and the proof are the article's own lines, in the article's own notation, and no retained line is substituted or re-flowed.  Retained uncounted as the source-local context the proof needs: lines 335--348.  Nothing was omitted from any retained span, so the package carries no omission record.  No retained line contains a citation, so the wrapper carries no bibliography and no bibliography entry.  No retained line contains a figure environment, an image-inclusion command or drawing code, so no image is delivered and the package declares no asset dependency.  Editorial formatting only, in addition: an AMS \texttt{amsart} preamble that restates the article's own declarations and macros used by the retained lines, namely the environments \texttt{lem}, \texttt{cor} on separate counters, together with the standard packages those lines need.  The retained environments print in this document as Lemma 1, Corollary 1 in order of appearance, whereas the article prints the same items with the numbering of its own full document.  The complete native source of the article as distributed in the arXiv e-print for this exact version is bundled unchanged as \texttt{sources/178429/source0.tex}.
\begin{lem} \label{magic_lem1}
Consider the mixture $\mu = p\mu_1 + (1-p)\mu_2$, where  $p \in (1/2, 1)$,  $\mu_1$ is a QID distribution with  triplet $(\gamma, \sigma^2, \nu)$, and $\mu_2$ is some distribution on \(\R.\) Suppose  that there exists a finite signed measure $\Lambda$ on $\mathbb{R}$ such that its Fourier transform is equal to 
\begin{eqnarray}\label{mu12}
\F\bigl[\Lambda \bigr](u) = \frac{\phi_2(u)}{\phi_1(u)}, \qquad u \in \R,
\end{eqnarray}
where   $\phi_1$ and $\phi_2$ are the characteristic functions of $\mu_1$ and $\mu_2$, If this measure is finite with total variation $|\Lambda|(\mathbb{R}) < p/(1-p)$, then $\mu$ is QID with the characteristic triplet \[\Bigl(\gamma + \int_{\mathbb{R}}c(x)\Tilde{\nu}(\mathrm{d}x),  \quad \sigma^2, \quad  \nu + \Tilde{\nu}\Bigr),\]
%where $c(x)$ is some representation function 
where $\Tilde{\nu}$ is a finite signed measure defined as follows: 
\begin{equation}\label{Tildenu}
    \Tilde{\nu}(B) = \sum_{m = 1}^{\infty} \frac{(-1)^{m + 1}}{m}\Bigl( 
\frac{1-p}{p} 
\Bigr)^m\Lambda^{*m}(B), \qquad B \in \B(\R\setminus\{0\}),
\end{equation}
with $\Lambda^{*m}(B)$ being the convolution of the measure $\Lambda$ with itself  $m$ times.\end{lem}

Note that the assumption $|\Lambda|(\mathbb{R}) < p/(1-p)$ is fulfilled if \(\Lambda\) is a probability measure, since \(p \in (1/2,1)\). The next two corollaries present specific cases when \(\Lambda\) is a probability measure, and therefore Lemma~\ref{magic_lem1} can be applied.\begin{cor}
(Cuppens theorem.) Any distribution \(\mu\) with an atom of mass larger than 1/2 is QID. 
\end{cor}

\begin{proof}
In fact, in this case \(\mu = p \mu_1 + (1-p) \mu_2,\) where \(p=\mu(\{\lambda\})>1/2\), \(\mu_1 = \delta_\lambda\), and 
\(\mu_2 = (\mu - p \delta_\lambda)/(1-p).\) Note that \(\mu_2\) is a probability measure, and \(\Lambda=\mu_2 * \delta_{-\lambda}\)   is also a probability measure. We conclude that   \(\mu\)  is a QID distribution with triplet \((\lambda,0, \nu),\)  where 
\begin{eqnarray*}
\nu = \sum_{m=1}^\infty \frac{(-1)^{m+1}}{m}
\Bigl( 
\frac{1-p}{p} 
\Bigr)^m \bigl(  \mu_2 * \delta_{-\lambda} \bigr)^{*m}.
\end{eqnarray*}
\end{proof}

\end{document}
